\documentclass[10pt,a4paper]{amsart}
\usepackage[margin=19mm]{geometry}
\usepackage{amsmath,amssymb,mathtools}
\usepackage[scr=rsfs,cal=euler]{mathalfa}
\usepackage{xfrac}
\usepackage[unicode,hidelinks,bookmarks=false]{hyperref}
\newcommand{\ie}{i.e.\ }
\newcommand{\cf}{cf.\ }
\newcommand{\resp}{respectively\ }
\newcommand{\Subsection}{\subsection}
\providecommand{\nobreakdash}{\nobreak\hbox{-}}
\setlength{\emergencystretch}{2em}
\multlinegap0pt
\usepackage{xcolor}
\usepackage{amssymb}
\usepackage{bm}
\usepackage[shortlabels]{enumitem}
\usepackage{tikz}
\usepackage{algorithm}\usepackage{algpseudocode}
\newtheorem{proposition}{Proposition}
\usepackage{cleveref}
\crefname{proposition}{Proposition}{Propositions}
\title{Optimizer-Induced Mode Connectivity: From AdamW to Muon}
\author{Fangzhao Zhang, Sungyoon Kim, Erica Zhang, Yiqi Jiang, Mert Pilanci}
\date{Source: arXiv:2605.09991v1 (CC BY 4.0)}
\begin{document}
\maketitle

\noindent\emph{Source and licence.} Optimizer-Induced Mode Connectivity: From AdamW to Muon. Version arXiv:2605.09991v1 (CC BY 4.0), distributed by arXiv under the Creative Commons Attribution 4.0 International licence, http://creativecommons.org/licenses/by/4.0/, as recorded in the arXiv licence metadata for this exact version and shown on the abstract page https://arxiv.org/abs/2605.09991v1. Full licence notice: \texttt{licenses/arxiv-2605.09991-CC-BY-4.0.txt}.\par\smallskip

\noindent\textit{Editorial scope.} Counted unit: the article's proposition prop5:computation (source0.tex lines 2132--2247). The article's complete original proof of it is delivered with it (source0.tex lines 2249--2622, one \texttt{proof} environment of 374 lines). No mathematics is written, repaired or re-derived: the statement and the proof are the article's own lines, in the article's own notation, and no retained line is substituted or re-flowed.  Retained uncounted as the source-local context the proof needs: lines 2072--2087.  Nothing was omitted from any retained span, so the package carries no omission record.  No retained line contains a citation, so the wrapper carries no bibliography and no bibliography entry.  No retained line contains a figure environment, an image-inclusion command or drawing code, so no image is delivered and the package declares no asset dependency.  Editorial formatting only, in addition: an AMS \texttt{amsart} preamble that restates the article's own declarations and macros used by the retained lines, namely the environments \texttt{proposition} on separate counters, together with the standard packages those lines need.  The retained environments print in this document as Proposition 1 in order of appearance, whereas the article prints the same items with the numbering of its own full document.  The complete native source of the article as distributed in the arXiv e-print for this exact version is bundled unchanged as \texttt{sources/200273/source0.tex}.
\begin{proposition}
\label{prop4:formula_Rop_Rinf}
Let $A\in\mathbb R^{d\times d}$ be invertible and
$y=\mathbf 1_{2d}$. Then for every $\sigma\in\{\pm1\}^d$,
\[
R_\infty(\mathcal C_{\sigma\mid-\sigma})
=
\|A^{-1}\sigma\|_\infty^{1/2},
\]
and
\[
R_{\mathrm{op}}(\mathcal C_{\sigma\mid-\sigma})
=
\sqrt{2\|A^{-1}\sigma\|_2}.
\]
\end{proposition}

\begin{proposition}
\label{prop5:computation}
Let $d\ge2$ and fix $1<L<\sqrt d$.
Denote
\[
h_1:=\mathbf 1=(1,\dots,1)^\top,
\qquad
h_2:=(1,-1,\dots,-1)^\top,
\]
and let
\[
B =
\begin{bmatrix}
\dfrac{1+L}{2}
&
\dfrac{1-L}{2(d-1)}
&
\dfrac{1-L}{2(d-1)}
&
\cdots
&
\dfrac{1-L}{2(d-1)}
\\[6pt]
\dfrac12
&
1-\dfrac{1}{2(d-1)}
&
-\dfrac{1}{2(d-1)}
&
\cdots
&
-\dfrac{1}{2(d-1)}
\\[6pt]
\dfrac12
&
-\dfrac{1}{2(d-1)}
&
1-\dfrac{1}{2(d-1)}
&
\cdots
&
-\dfrac{1}{2(d-1)}
\\
\vdots
&
\vdots
&
\vdots
&
\ddots
&
\vdots
\\[4pt]
\dfrac12
&
-\dfrac{1}{2(d-1)}
&
-\dfrac{1}{2(d-1)}
&
\cdots
&
1-\dfrac{1}{2(d-1)}
\end{bmatrix}.
\]
Set
\[
A:=B^{-1},
\qquad
y=\mathbf 1_{2d}.
\]
Then $R_\infty(\mathcal C_{\sigma\mid-\sigma})$
is minimized exactly at $\sigma=\pm h_1$, while
$
R_{\mathrm{op}}(\mathcal C_{\sigma\mid-\sigma})
$
is minimized exactly at $\sigma=\pm h_2$. Moreover, 
Assume $L=\frac{\sqrt d}{2}\gg 2$. Then for sufficiently large $d \geq 16$,
\[
\min_{\sigma\in\{\pm1\}^d}
R_\infty(\mathcal C_{\sigma\mid-\sigma})
=
R_\infty(\mathcal C_{h_1\mid-h_1})
=
1, \quad 
\min_{\substack{\sigma\in\{\pm1\}^d\\ \sigma\neq \pm h_1}}
R_\infty(\mathcal C_{\sigma\mid-\sigma})
=
\left(
1+\frac{\sqrt d/2-1}{d-1}
\right)^{1/2},
\]
\[
\min_{\sigma\in\{\pm1\}^d}
R_{\mathrm{op}}(\mathcal C_{\sigma\mid-\sigma})
=
R_{\mathrm{op}}(\mathcal C_{h_2\mid-h_2})
=
d^{1/4}/\sqrt{2},
\]
and
\[
\min_{\substack{\sigma\in\{\pm1\}^d\\ \sigma\neq \pm h_2}}
R_{\mathrm{op}}(\mathcal C_{\sigma\mid-\sigma})
=
\sqrt2
\left[
\left(
\frac{\sqrt d}{2}
-
\frac{\frac{\sqrt d}{2}-1}{d-1}
\right)^2
+
4-\frac{3}{d-1}
\right]^{1/4}.
\]
\end{proposition}

\begin{proof}
By \cref{prop4:formula_Rop_Rinf},
\[
R_\infty(\mathcal C_{\sigma\mid-\sigma})
=
\|B\sigma\|_\infty^{1/2},
\qquad
R_{\mathrm{op}}(\mathcal C_{\sigma\mid-\sigma})
=
\sqrt{2\|B\sigma\|_2}.
\]
Hence it suffices to minimize $\|B\sigma\|_\infty$ and $\|B\sigma\|_2$ over
$\sigma\in\{\pm1\}^d$.
Both norms are invariant under $\sigma\mapsto-\sigma$, since
$B(-\sigma)=-B\sigma$. Thus we may assume $\sigma_1=1$.

Let
\[
k:=|\{i\in\{2,\dots,d\}:\sigma_i=-1\}|,
\qquad
r:=\frac{k}{d-1}.
\]
Then $0\le r\le1$. From the definition of $B$,
\[
(B\sigma)_1=1+(L-1)r,
\]
and for $i\ge2$,
\[
(B\sigma)_i=
\begin{cases}
1+r, & \sigma_i=1,\\
-1+r, & \sigma_i=-1.
\end{cases}
\]

We first minimize $\|B\sigma\|_\infty$. If $r=0$, then $\sigma=h_1$ and
\[
B\sigma=\mathbf 1,
\]
so $\|B\sigma\|_\infty=1$. If $r>0$, then since $L>1$,
\[
|(B\sigma)_1|=1+(L-1)r>1.
\]
Thus the unique minimizer with $\sigma_1=1$ is $\sigma=h_1$. By sign symmetry,
the global minimizers are exactly $\sigma=\pm h_1$.

Now minimize $\|B\sigma\|_2$. The coordinate formula gives
\[
\|B\sigma\|_2^2
=
\bigl(1+(L-1)r\bigr)^2
+
(d-k-1)(1+r)^2
+
k(-1+r)^2.
\]
Using $k=(d-1)r$, the last two terms become
\[
(d-1)(1-r)(1+r)^2+(d-1)r(1-r)^2.
\]
Since
\[
(1-r)(1+r)^2+r(1-r)^2
=
1+2r-3r^2,
\]
we obtain
\[
\|B\sigma\|_2^2
=
\bigl(1+(L-1)r\bigr)^2
+
(d-1)(1+2r-3r^2).
\]
For $\sigma=h_2$, we have $r=1$, hence
\[
Bh_2=Le_1,
\qquad
\|Bh_2\|_2^2=L^2.
\]

For $r<1$, set $\delta:=1-r>0$. Then
\[
1+(L-1)r
=
L-(L-1)\delta,
\]
and
\[
1+2r-3r^2
=
4\delta-3\delta^2.
\]
Therefore
\[
\|B\sigma\|_2^2
=
\bigl(L-(L-1)\delta\bigr)^2
+
(d-1)(4\delta-3\delta^2).
\]
Subtracting $L^2$ gives
\[
\begin{aligned}
\|B\sigma\|_2^2-L^2
&=
-2L(L-1)\delta
+
(L-1)^2\delta^2
+
4(d-1)\delta
-
3(d-1)\delta^2
\\
&=
\delta
\left[
4(d-1)-2L(L-1)
+
\delta\bigl((L-1)^2-3(d-1)\bigr)
\right].
\end{aligned}
\]
We claim the bracket is strictly positive for every $\delta\in(0,1]$.

If
\[
(L-1)^2-3(d-1)\ge0,
\]
then the bracket is minimized at $\delta=0$, where it equals
\[
4(d-1)-2L(L-1).
\]
This is strictly positive because
\[
4(d-1)-2L(L-1)
>
4(d-1)-2d
=
2d-4
\ge0,
\]
and the inequality is strict since $L(L-1)<L^2<d$.

If
\[
(L-1)^2-3(d-1)<0,
\]
then the bracket is minimized at $\delta=1$, where it equals
\[
\begin{aligned}
&4(d-1)-2L(L-1)+(L-1)^2-3(d-1)
\\
&=
d-L^2
>
0.
\end{aligned}
\]
Therefore
\[
\|B\sigma\|_2^2>L^2
\]
whenever $r<1$. Thus the unique minimizer with $\sigma_1=1$ is $\sigma=h_2$. By sign symmetry,
the global minimizers are exactly $\sigma=\pm h_2$.

We now derive the runner-up formulas. For the $R_\infty$ runner-up, recall that
\[
\|B\sigma\|_\infty
=
\max\{1+(L-1)r,\ 1+r\},
\qquad
r=\frac{k}{d-1}.
\]
The minimum occurs at $r=0$. The smallest positive value of $r$ is
\[
r=\frac1{d-1}.
\]
Therefore
\[
\left(R_\infty^{(2)}\right)^2
=
\max\left\{
1+\frac{L-1}{d-1},
1+\frac1{d-1}
\right\}
=
1+\frac{\max\{L-1,1\}}{d-1}.
\]
Hence
\[
R_\infty^{(2)}
=
\left(
1+\frac{\max\{L-1,1\}}{d-1}
\right)^{1/2}.
\]

For the $R_{\mathrm{op}}$ runner-up, define
\[
F(r):=\|B\sigma\|_2^2
=
\bigl(1+(L-1)r\bigr)^2
+
(d-1)(1+2r-3r^2).
\]
Then
\[
R_{\mathrm{op}}(\mathcal C_{\sigma\mid-\sigma})
=
\sqrt2\,F(r)^{1/4}.
\]
The minimum occurs at $r=1$, where $F(1)=L^2$.
Thus the runner-up is the smallest value of $F(r)$ over
\[
r\in
\left\{
0,\frac1{d-1},\dots,1-\frac1{d-1}
\right\}.
\]
Equivalently, set $\delta:=1-r$. Then
\[
\delta\in
\left\{
\frac1{d-1},\frac2{d-1},\dots,1
\right\}.
\]
From the previous computation,
\[
F(r)-L^2
=
\delta
\left[
4(d-1)-2L(L-1)
+
\delta\bigl((L-1)^2-3(d-1)\bigr)
\right].
\]
As a function of $\delta$, this is a quadratic with leading coefficient
\[
(L-1)^2-3(d-1).
\]
Since $L^2<d$, we have
\[
(L-1)^2 < L^2 < d < 3(d-1)
\]
for $d\ge2$, and hence the quadratic is strictly concave. Therefore its minimum over the discrete interval occurs at one of the endpoints:
\[
\delta=\frac1{d-1}
\qquad\text{or}\qquad
\delta=1.
\]
Equivalently,
\[
r=1-\frac1{d-1}
\qquad\text{or}\qquad
r=0.
\]
Hence
\[
F^{(2)}
=
\min\left\{
F\left(1-\frac1{d-1}\right),
F(0)
\right\}.
\]
Now
\[
F(0)=d,
\]
and
\[
\begin{aligned}
F\left(1-\frac1{d-1}\right)
&=
\left(
1+(L-1)\left(1-\frac1{d-1}\right)
\right)^2
\\
&\quad+
(d-1)
\left[
1+2\left(1-\frac1{d-1}\right)
-
3\left(1-\frac1{d-1}\right)^2
\right]
\\
&=
\left(
L-\frac{L-1}{d-1}
\right)^2
+
4-\frac{3}{d-1}.
\end{aligned}
\]
Therefore
\[
F^{(2)}
=
\min\left\{
d,\,
\left(
L-\frac{L-1}{d-1}
\right)^2
+
4-\frac{3}{d-1}
\right\}.
\]
Since
\[
R_{\mathrm{op}}=\sqrt2\,F(r)^{1/4},
\]
we get
\[
R_{\mathrm{op}}^{(2)}
=
\sqrt2
\left[
\min\left\{
d,\,
\left(
L-\frac{L-1}{d-1}
\right)^2
+
4-\frac{3}{d-1}
\right\}
\right]^{1/4}.
\]

Finally, set $L=\sqrt d/2$. If $d\ge16$, then $L-1\ge1$, and hence
\[
R_\infty^{(2)}
=
\left(
1+\frac{\frac{\sqrt d}{2}-1}{d-1}
\right)^{1/2}.
\]
Moreover,
\[
R_{\mathrm{op}}(\mathcal C_{h_2\mid-h_2})
=
\sqrt{2L}
=
d^{1/4}.
\]
For sufficiently large $d$, one also has
\[
\left(
\frac{\sqrt d}{2}
-
\frac{\frac{\sqrt d}{2}-1}{d-1}
\right)^2
+
4-\frac{3}{d-1}
< d,
\]
so the minimum in the formula for $F^{(2)}$ is attained by the second term. Therefore
\[
R_{\mathrm{op}}^{(2)}
=
\sqrt2
\left[
\left(
\frac{\sqrt d}{2}
-
\frac{\frac{\sqrt d}{2}-1}{d-1}
\right)^2
+
4-\frac{3}{d-1}
\right]^{1/4}.
\]
This proves the claimed formulas.
\end{proof}

\end{document}
